\section{A uniform preconditioner that retains arithmetic (NS-80)}
\label{sec:ns80-arithmetic-preconditioner}

The complete sampling comparison also bounds every finite Gram matrix.
Its comparator retains the arithmetic values $A(m/n)$ at \emph{all}
integer knots; it is not the smooth model from NS-70. This gives a
uniform finite-block efficiency theorem, but neither a fast inverse
nor the cofinal lower bound on available gain needed for RH.

Define the complete arithmetic sampling Gram
\begin{equation}
 Q_{ij}=\frac1{ij}+\sum_{m=1}^\infty\frac{
 A(m/i)A(m/j)+A((m+1)/i)A((m+1)/j)}{m(m+1)}.
 \label{eq:ns80-sampling-gram}
\end{equation}
The series converges by NS-78 and Cauchy--Schwarz. For any finite
coefficient vector, its quadratic form is exactly the complete
sampling quantity for the corresponding pure atom combination.
Thus on every finite index set,
\begin{equation}
 \frac1{16}Q\preceq G\preceq\frac{21}{16}Q.
 \label{eq:ns80-full-comparison}
\end{equation}
Both matrices are positive definite: zero physical norm forces zero
samples and exterior coefficient, and the exact inversion in NS-78
then forces every coefficient to vanish.

\begin{proposition}[Uniform efficiency with the full arithmetic comparator]
\label{prop:ns80-uniform-efficiency}
Use the ordered basis consisting of the old $N$ atoms and the new
adjacent differences. Let $H_G,H_Q$ be the respective Schur complements
of the old blocks in their complete Grams. Then
\[
 \frac1{16}H_Q\preceq H_G\preceq\frac{21}{16}H_Q.
\]
For any nonzero \emph{actual} residual correlation vector $g$, set
$v=H_Q^{-1}g$. Its optimally scaled true gain satisfies
\begin{equation}
 \frac{(g^Tv)^2}{v^TH_Gv}
 \geq\frac{21}{121}\,g^TH_G^{-1}g.
 \label{eq:ns80-gain-fraction}
\end{equation}
This is uniform in $N$ and $g$. All old projections and all cross terms
are included. It is a fraction of the full available block gain, not
of the current residual energy.
\end{proposition}
\begin{proof}
The full comparison survives the invertible change to adjacent
coordinates. For each new vector $v$, the Schur quadratic form is the
minimum of the complete quadratic form over old coordinates. Taking
those minima on both sides preserves the same two constants; the old
minimizers need not agree.

More generally put $a=1/16$, $b=21/16$ and
$M=H_Q^{-1/2}H_GH_Q^{-1/2}$. Its eigenvalues lie in $[a,b]$.
With $s=H_Q^{-1/2}g$, normalize its squared spectral coordinates to
probability weights and write $x=\mathbb E\lambda$,
$y=\mathbb E\lambda^{-1}$. For $a\leq\lambda\leq b$,
\[
 \lambda+\frac{ab}{\lambda}\leq a+b,
 \qquad x+aby\leq a+b,
 \qquad xy\leq\frac{(a+b)^2}{4ab}.
\]
The final inequality follows from $(x-aby)^2\geq0$.
The ratio of the selected gain to the full gain is precisely
$1/(xy)$, hence is at least $4ab/(a+b)^2=21/121$.
This is the usual elementary Kantorovich argument, displayed here
to make every hypothesis and normalization explicit.
\end{proof}

\paragraph{Repeated steps within a fixed finite block.}
Start at $x_0=0$ and minimize the exact block quadratic
$\Phi(x)=x^TH_Gx-2g^Tx$. At each step use the current residual
$r_k=g-H_Gx_k$, direction $p_k=H_Q^{-1}r_k$, and the exact line search
\[
 x_{k+1}=x_k+\frac{r_k^Tp_k}{p_k^TH_Gp_k}p_k,
\]
with termination if $r_k=0$. The remaining block gain is
$r_k^TH_G^{-1}r_k$, so Proposition~\ref{prop:ns80-uniform-efficiency}
reduces it by at least $21/121$ each time. After $k$ steps the captured
fraction of full gain is therefore at least
\[
 1-(100/121)^k.
\]
In particular twenty ideal steps capture more than $97.79\%$, uniformly
in block size. This statement assumes application of the complete
$H_Q^{-1}$ and the true $H_G$; no implementation or complexity bound for
that infinite-sample comparator is asserted. Truncating its sum without
a certified tail would invalidate the theorem's premises.

For completeness, minimizing the full sampling objective over the first
$N$ coefficients gives a residual with squared error at most $21E_N$:
apply both bounds of NS-78 before and after the two minimizations.
This uniform approximation factor is also not a decay estimate for $E_N$.

\paragraph{The remaining obstruction.}
Even perfect capture of every block gain does not prove that the limiting
error is zero. The full block gains telescope to $E_N-E_\infty$, and
NS-76 exhibits a component invisible to all of them. A lower estimate
for available gain relative to the entire current error remains essential.
The present theorem preserves arithmetic in its comparator and supplies a
uniform finite-block efficiency guarantee; it does not provide that lower
estimate, a fast computable replacement for $H_Q$, a bound for the earlier
smooth curvature rule, G2 or RH.
